\section{Introduction}
\label{sec:intro}

Universal algebra studies classes of algebraic structures defined by equational
laws: identities of the form $s = t$ where $s$ and $t$ are terms built from
variables and operation symbols.  A fundamental question is the \emph{implication
problem}: given equational laws $E$ and $E'$, does $E$ logically entail $E'$?
That is, does every model of $E$ also satisfy $E'$?

The \emph{Equational Theories Project} (ETP) \cite{etp2025} recently determined
the complete implication relation among $4{,}694$ equational laws over a single
binary operation $*$ (a \emph{magma}) of \emph{order at most four} (at most four
applications of $*$), verifying each of the $22{,}028{,}942$ pairwise implication
decisions with formal Lean~4 proofs.  This dataset provides a unique benchmark:
a complete, machine-verified ground truth against which any automated implication
classifier can be tested.

\para{The classification problem.}
A practical implication classifier must be both \emph{sound} (it never asserts
a false implication) and \emph{complete} (it never misses a true implication or
counterexample).  Full decidability of the implication problem for unrestricted
magma laws is known to be unattainable \cite{mckenzie1971spectra}; nonetheless,
efficient heuristic and semi-decision procedures are valuable for interactive
algebraic exploration and automated hypothesis testing.

\para{Soundness gaps in prior classifiers.}
The predecessor classifier \textsc{predictor3} employs two components that are
either unsound or incomplete.  First, a bounded breadth-first reachability search
serves as its positive implication test.  While sound in principle---a reachable
term intersection genuinely witnesses an equational derivation---its bounded
character means it can spuriously report an implication when deep rewrite chains
cause the reachable-term sets to collide accidentally, without a true proof.
Second, a heuristic hill-climbing search tests for counterexamples in magmas of
sizes 4 and 5; because hill-climbing is not exhaustive, it may terminate without
finding a genuine counterexample of those sizes.

\para{Our approach.}
We replace both unsound or incomplete components with established sound and
complete procedures from the term rewriting and satisfiability literatures.

For the \emph{positive path}, we implement Knuth-Bendix (KB) completion
\cite{knuth1970simple} from scratch in Python.  When KB completion terminates
with a convergent term rewriting system $R$, the normal-form equality test
$\mathrm{NF}_R(\ell_2) = \mathrm{NF}_R(r_2)$ is a sound proof of $E \models E'$
by the Knuth-Bendix soundness theorem and Birkhoff's completeness theorem
\cite{birkhoff1935structure}.  The bounded BFS search serves as a fallback for
the approximately 88\% of cases where KB does not terminate within the allotted time.

For the \emph{negative path}, we encode the finite model finding problem as an
SMT query in Z3 \cite{demoura2008z3} for domain sizes $n = 4$ and $n = 5$.
Z3 is complete for each fixed finite domain: if no counterexample of size~$n$
exists, Z3 returns \textsc{unsat}, providing stronger evidence for implication
than a heuristic search that merely fails to find a witness.

\para{The soundness invariant.}
During evaluation, we discovered that a buggy KB implementation generated rewrite
rules $\ell \to r$ with $\operatorname{vars}(r) \not\subseteq \operatorname{vars}(\ell)$---specifically,
rules of the form $x \to y$ where $x$ and $y$ are distinct variables.  Such rules
are not valid term rewriting rules under any reduction order, yet they can arise
from incorrect variable renaming in the critical pair computation.  We prove
(Theorem~\ref{thm:collapse}) that any TRS containing a rule with a free RHS variable
collapses all normal forms to a single term, making the normal-form test trivially
true for all equation pairs and producing an unlimited supply of false positives.
Enforcing the invariant $\operatorname{vars}(r) \subseteq \operatorname{vars}(\ell)$ before accepting
any rule eliminates this failure mode.

\para{Main contributions.}
\begin{itemize}[leftmargin=*,itemsep=2pt,topsep=2pt]
  \item We formalize the \emph{soundness invariant} for KB completion in the equational
        implication setting and prove that its violation causes total normal-form collapse
        (Theorem~\ref{thm:collapse}).
  \item We prove that, under the soundness invariant, KB completion provides a genuine
        proof of equational implication with no false positives
        (Theorem~\ref{thm:kb-soundness}).
  \item We establish Z3's completeness for finite magma counterexample search at fixed
        domain sizes (Proposition~\ref{prop:z3-complete}), replacing heuristic hill-climbing
        with a decision procedure.
  \item We show that the resulting layered classifier \textsc{predictor4} produces no false
        positives from its KB component and achieves $100\%$ accuracy on stratified samples
        of 50 and 200 implications from the ETP ground truth
        (Theorem~\ref{thm:no-fp} and Section~\ref{sec:experiments}).
\end{itemize}

\para{Organization.}
Section~\ref{sec:prelim} establishes notation and the necessary background on magmas,
equational implication, and term rewriting.  Section~\ref{sec:main} presents the main
theoretical results.  Section~\ref{sec:classifier} describes the layered decision
procedure \textsc{predictor4}.  Section~\ref{sec:experiments} reports the empirical
evaluation.  Section~\ref{sec:discussion} discusses limitations and open questions, and
Section~\ref{sec:conclusion} concludes.
